% ============================================================
% Paper deep reading template
% File name suggestion: YYYY_MM_DD_short_name.tex
% ============================================================

\section{论文精读：Paper Title}

\begin{conceptbox}
\textbf{论文：} \href{https://arxiv.org/abs/0000.00000}{Paper Title}\\
\textbf{作者/机构：} TODO\\
\textbf{时间/venue：} TODO\\
\textbf{方向标签：} TODO，例如 dLLM / post-training / image editing / world model\\
\textbf{阅读日期：} TODO
\end{conceptbox}

\subsection{0. 一句话结论}

TODO：用一两句话说清楚这篇论文真正贡献了什么，而不是复述标题。

\begin{summarybox}
\textbf{我的判断：} TODO：值得精读 / 只需速读 / 作为背景引用 / 需要等待代码验证。
\end{summarybox}

\subsection{1. 为什么读这篇}

\begin{itemize}[leftmargin=1.5em]
  \item \textbf{和当前主线的关系：} TODO。
  \item \textbf{它可能补上的知识缺口：} TODO。
  \item \textbf{我希望读完回答的问题：} TODO。
\end{itemize}

\subsection{2. 问题定义与背景}

\textbf{问题：} TODO：这篇论文解决的具体问题是什么？

\textbf{输入/输出：} TODO：模型或算法的输入、输出、训练数据、评价对象分别是什么？

\textbf{前作限制：} TODO：它认为已有方法哪里不够好？

\begin{intubox}
\textbf{直觉版：} TODO：不用公式，用自己的话讲这篇论文的核心想法。
\end{intubox}

\subsection{3. 方法拆解}

\subsubsection{3.1 总体框架}

TODO：画出或文字描述 pipeline。可以回答：

\begin{itemize}[leftmargin=1.5em]
  \item 模型结构有什么变化？
  \item 训练目标有什么变化？
  \item 数据构造或采样流程有什么变化？
  \item 推理阶段有什么特殊设计？
\end{itemize}

\subsubsection{3.2 关键机制}

\begin{itemize}[leftmargin=1.5em]
  \item \textbf{机制一：} TODO。
  \item \textbf{机制二：} TODO。
  \item \textbf{机制三：} TODO。
\end{itemize}

\subsubsection{3.3 关键公式或算法}

TODO：只抄真正需要的公式，并写清楚每个符号的含义。

\begin{warnbox}
\textbf{注意：} 不要只复制公式。这里要写“这个公式为什么这样设计，以及如果去掉某一项会怎样”。
\end{warnbox}

\subsection{4. 实验与证据}

\begin{itemize}[leftmargin=1.5em]
  \item \textbf{主要 benchmark：} TODO。
  \item \textbf{核心对比对象：} TODO。
  \item \textbf{最有信息量的 ablation：} TODO。
  \item \textbf{结论是否被实验充分支撑：} TODO。
\end{itemize}

\subsection{5. 局限与风险}

\begin{itemize}[leftmargin=1.5em]
  \item \textbf{假设限制：} TODO。
  \item \textbf{实验限制：} TODO。
  \item \textbf{工程限制：} TODO。
  \item \textbf{可能失败场景：} TODO。
\end{itemize}

\subsection{6. 与已有工作的关系}

\begin{itemize}[leftmargin=1.5em]
  \item \textbf{直接前作：} TODO。
  \item \textbf{相近路线：} TODO。
  \item \textbf{真正差异：} TODO。
\end{itemize}

\subsection{7. 对我的可复用点}

\begin{itemize}[leftmargin=1.5em]
  \item \textbf{可复用 idea：} TODO。
  \item \textbf{可复用实验设计：} TODO。
  \item \textbf{可复用写法/图表：} TODO。
  \item \textbf{值得追的 follow-up：} TODO。
\end{itemize}

\subsection{8. 待查问题}

\begin{itemize}[leftmargin=1.5em]
  \item TODO。
  \item TODO。
  \item TODO。
\end{itemize}

